\documentclass[aps,prd,preprint,amsmath,amssymb,nofootinbib]{revtex4-2}

\usepackage{graphicx}
\usepackage{bm}
\usepackage{hyperref}
\hypersetup{colorlinks=true,linkcolor=blue,citecolor=blue,urlcolor=blue}

\begin{document}

\title{A precise numerical relation linking the tau--electron mass ratio
to the electromagnetic and gravitational coupling strengths}

\author{Thomas J.\ Buckholtz}
\affiliation{National Coalition of Independent Scholars, Brattleboro, Vermont 05301, USA}
\affiliation{Ronin Institute for Independent Scholarship, Montclair, New Jersey 07043, USA}
\email{tjb@alumni.caltech.edu}

\date{\today}

\begin{abstract}
We report a numerical relation among measured quantities,
$(4/3)\,(m_\tau/m_e)^{12} = \alpha_{\mathrm{EM}}/\alpha_G$,
in which $\alpha_{\mathrm{EM}}$ is the fine-structure constant and
$\alpha_G \equiv G m_e^2/(\hbar c)$ is the gravitational coupling strength
referred to the electron mass. Using current Particle Data Group and CODATA
inputs, the two sides agree to $0.0135\%$, which corresponds to
$0.17\,\sigma$ of the present tau-mass uncertainty. To assess whether such
agreement is likely to arise by chance, we search all relations of the form
$(p/q)(m_i/m_j)^n$ with $p,q\le 10$, $1\le n\le 24$, and $m_i,m_j$ drawn from
the eleven charged fermion, proton and neutron masses; among the resulting
$83\,160$ candidates the present relation ranks first, the next-best candidate
being a factor of $22$ less accurate. The empirical probability of a chance
match at the one-percent level is below $10^{-4}$. We present the relation,
its statistical significance, the algebraic structure of its integer
coefficients, and the questions it leaves open---foremost the derivation of
the coefficient $4/3$, which we find is not supplied by any adjacent
theoretical tradition. We suggest the relation belongs, in a static
(non-time-varying) form, to the tradition of large-number coincidences that
involve the gravitational coupling.
\end{abstract}

\maketitle

\section{Introduction}

Empirical relations among the masses of the elementary particles, and between
those masses and the fundamental coupling strengths, have long been used as
signposts toward a deeper theory~\cite{koide1983,dirac1938,wesson1980}. The
best known such relation among the charged leptons is Koide's
formula~\cite{koide1983}, which involves the three charged-lepton masses alone
and makes no reference to gravitation. Separately, P.~A.~M.~Dirac's
large-numbers hypothesis~\cite{dirac1938} drew attention to the enormous
dimensionless ratio between the electromagnetic and gravitational coupling
strengths, and to its apparent coincidence with cosmological quantities.

Here we report a relation that connects these two themes: it links a
charged-lepton mass ratio to the ratio of the electromagnetic and
gravitational coupling strengths. We state the relation, quantify how unlikely
it is to be an accident, examine the arithmetic of its coefficients, and---in
the interest of full disclosure---set out plainly what the relation does and
does not establish. Our intent is to place a single, checkable numerical fact
on the record, not to claim a mechanism for it.

\section{The relation}

We define the gravitational coupling strength referred to the electron mass,
in the manner introduced by Dirac~\cite{dirac1938},
\begin{equation}
\alpha_G \equiv \frac{G\,m_e^2}{\hbar c},
\label{eq:alphaG}
\end{equation}
which is the gravitational analogue of the fine-structure constant
$\alpha_{\mathrm{EM}} = e^2/(4\pi\varepsilon_0\hbar c)$. The choice of the
electron mass as the reference is a convention; we return to it in
Sec.~\ref{sec:discussion}.

The relation we report is
\begin{equation}
\frac{4}{3}\left(\frac{m_\tau}{m_e}\right)^{\!12}
= \frac{\alpha_{\mathrm{EM}}}{\alpha_G}.
\label{eq:main}
\end{equation}
With the inputs listed in Table~\ref{tab:inputs}, the two sides evaluate to
\begin{align}
\text{LHS} &= 4.16617\times 10^{42}, \nonumber\\
\text{RHS} &= 4.16561\times 10^{42},
\end{align}
so that
\begin{equation}
\frac{\text{LHS}}{\text{RHS}} = 1.000135,
\end{equation}
a fractional discrepancy of $0.0135\%$. Because the left-hand side scales as
$m_\tau^{12}$, the present tau-mass uncertainty
$\sigma_{m_\tau}=0.12~\mathrm{MeV}$ propagates to a one-standard-deviation band
of $12\,\sigma_{m_\tau}/m_\tau = 0.081\%$ on the left-hand side. The observed
discrepancy is therefore $0.17\,\sigma$. Equivalently, Eq.~(\ref{eq:main})
holds exactly for
\begin{equation}
m_\tau = 1776.840~\mathrm{MeV},
\end{equation}
which lies within the current experimental window
$1776.86\pm 0.12~\mathrm{MeV}$~\cite{pdg2024}. A measurement of the tau mass at
the $0.01~\mathrm{MeV}$ precision anticipated at Belle~II~\cite{belleii} would
sharpen this test by an order of magnitude.

\begin{table}[t]
\caption{Input quantities and their sources. Masses are in
$\mathrm{MeV}/c^2$.}
\label{tab:inputs}
\begin{ruledtabular}
\begin{tabular}{lll}
Quantity & Value & Source \\
\hline
$m_e$ & $0.51099895000$ & CODATA~2018~\cite{codata2018} \\
$m_\tau$ & $1776.86\pm 0.12$ & PDG~2024~\cite{pdg2024} \\
$\alpha_{\mathrm{EM}}$ & $1/137.035999084$ & CODATA~2018~\cite{codata2018} \\
$G$ & $6.67430\times 10^{-11}$ & CODATA~2018~\cite{codata2018} \\
$\alpha_G$ & $1.75181\times 10^{-45}$ & from Eq.~(\ref{eq:alphaG}) \\
\end{tabular}
\end{ruledtabular}
\end{table}

\section{Statistical significance}
\label{sec:lookelsewhere}

A single close numerical coincidence is not, by itself, evidence of anything:
with enough freedom in choosing masses, exponents and rational prefactors, one
can approximate almost any target. The relevant question is therefore a
look-elsewhere question---how special is Eq.~(\ref{eq:main}) within the family
of relations one might have written down instead?

To answer it we enumerate every relation of the form
\begin{equation}
\frac{p}{q}\left(\frac{m_i}{m_j}\right)^{\!n},
\qquad
1\le p,q \le 10,\quad
1\le n \le 24,\quad
i\neq j,
\label{eq:family}
\end{equation}
with $m_i,m_j$ drawn from the set
$\{e,\mu,\tau,u,d,s,c,b,t,p,n\}$, and compare each to the target
$\alpha_{\mathrm{EM}}/\alpha_G$. This defines a search space of $83\,160$
candidates. The results are summarized in Table~\ref{tab:lookelsewhere}.

\begin{table}[t]
\caption{Look-elsewhere audit of Eq.~(\ref{eq:main}) over the
$83\,160$-member family~(\ref{eq:family}). ``Hits'' counts candidates matching
the target to better than the stated tolerance.}
\label{tab:lookelsewhere}
\begin{ruledtabular}
\begin{tabular}{lr}
Quantity & Value \\
\hline
Total candidates & $83\,160$ \\
Hits within $5\%$ & $26$ \\
Hits within $1\%$ & $5$ \\
Hits within $0.1\%$ & $1$ \\
Rank of Eq.~(\ref{eq:main}) & $1$ \\
Error of Eq.~(\ref{eq:main}) & $0.0135\%$ \\
Second-best candidate & $(7/5)(m_t/m_s)^{13}$, $0.30\%$ \\
Empirical $p$ (within $1\%$) & $<10^{-4}$ \\
\end{tabular}
\end{ruledtabular}
\end{table}

Eq.~(\ref{eq:main}) is the unique candidate that matches the target to better
than $0.1\%$; it ranks first overall, and the next-best relation,
$(7/5)(m_t/m_s)^{13}$, is a factor of $22$ less accurate. Only five of the
$83\,160$ candidates match to better than $1\%$, giving an empirical
probability below $10^{-4}$ that a match of this quality would arise by chance
within this family. We emphasize that this analysis quantifies non-randomness;
it does not by itself supply a mechanism.

\subsection{Robustness to the choice of reference mass}
\label{sec:refmass}

A further concern is that $m_e$ was chosen as the reference mass in
$\alpha_G$ by convention (Sec.~\ref{sec:discussion}); one may ask whether an
equally precise relation would have emerged had a different reference been
chosen, in which case the apparent precision of Eq.~(\ref{eq:main}) would owe
more to the size of the search space than to the electron. We therefore
repeat the search of Eq.~(\ref{eq:family}) with the reference mass in
$\alpha_G$ itself allowed to vary over all eleven masses, pooling the
resulting $132\,000$ candidates into a single ranked list. Two of the
eleven reference choices produce a match better than $0.1\%$: the original
relation, and
\begin{equation}
\frac{4}{7}\left(\frac{m_t}{m_\tau}\right)^{\!18}
= \frac{\alpha_{\mathrm{EM}}}{\alpha_G(m_\tau)},
\label{eq:tau-alt}
\end{equation}
which matches its target to $0.0070\%$ using the top-quark--tau ratio---a
factor of two tighter than Eq.~(\ref{eq:main}). Eq.~(\ref{eq:tau-alt}) is not
a rescaling of Eq.~(\ref{eq:main}): it involves a different mass pair,
reference mass, and coefficient. Honestly reported, Eq.~(\ref{eq:main}) is
therefore the \emph{second}-best-fitting relation in the pooled sample of
$132\,000$ candidates, beaten by this one genuinely independent competitor.
Its pooled empirical probability---the fraction of candidates at least as
accurate, $3/132\,000$---is $2.3\times 10^{-5}$, comparable to (and, taken at
face value, marginally tighter than) the single-reference estimate of
Sec.~\ref{sec:lookelsewhere}; the family therefore remains a two-standout
outlier rather than a single one when the reference-mass convention is
relaxed. We record Eq.~(\ref{eq:tau-alt}) as an independent numerical
curiosity in its own right, not as a derivation of, or an argument against,
Eq.~(\ref{eq:main}); we take no position on whether it reflects anything
beyond coincidence.

\section{Algebraic structure of the coefficients}

The integers that appear in Eq.~(\ref{eq:main})---the prefactor $4/3$ and the
exponent $12$---admit suggestive algebraic readings, which we record as
observations, not as derivations.

The exponent $12$ coincides with two independent invariants of the exceptional
Lie algebras: it is the number of nonzero roots of
$\mathrm{G}_2 = \mathrm{Aut}(\mathbb{O})$, the automorphism group of the
octonions, and it is the highest of the four Casimir degrees $(2,6,8,12)$ of
$\mathrm{F}_4$~\cite{bincer1993}. The prefactor may be written
$4/3 = 36/27$, where $27 = \dim J_3(\mathbb{O})$ is the dimension of the
exceptional (Albert) Jordan algebra and $36 = \dim \mathrm{SO}(9)$. A search
over rational combinations of dimensions of classical and exceptional Lie
algebras finds both the ratio $4/3$ and the specific triple
$\{36,27,12\}$ with an estimated chance probability of order $1/37000$.

We flag this algebraic search explicitly as a \emph{secondary}, post-hoc
look-elsewhere test: unlike the primary search of
Sec.~\ref{sec:lookelsewhere}, which was defined and run without reference to
the value $4/3$, the algebraic search was performed \emph{after} $4/3$ was
already known, over a pool of dimensions chosen because they were expected to
contain it. Its $1/37000$ figure should therefore be read as a measure of how
special the triple $\{36,27,12\}$ is \emph{within that pool}, not as an
independent confirmation of Eq.~(\ref{eq:main}) on the same footing as
Sec.~\ref{sec:lookelsewhere}.

These connections identify the algebraic objects that \emph{count} the
coefficients, but they do not explain \emph{why} those objects should govern
the ratio $\alpha_{\mathrm{EM}}/\alpha_G$. In particular we have been unable to
derive the coefficient $4/3$ from first principles, and---as discussed in
Sec.~\ref{sec:discussion}---we have not found such a derivation in any adjacent
theoretical tradition.

\section{Relation to prior work}

Eq.~(\ref{eq:main}) is distinct from the charged-lepton relations of the
Koide--Sumino type~\cite{koide1983,sumino2009} and from the exceptional-Jordan
mass-ratio programme~\cite{singh2021}. Those constructions predict ratios among
the fermion masses from flavour symmetry or algebra alone, and involve no
gravitational quantity; the present relation, by contrast, is characterized
precisely by the appearance of the gravitational coupling $\alpha_G$. We have
verified that the integer $12$ enters the algebra of Ref.~\cite{singh2021} as a
count of eigenvalues rather than as an exponent, so that the numerical
agreement of the two occurrences does not amount to a shared derivation.

The natural home for a relation between a particle-mass combination and
$\alpha_G$ is instead the tradition of large-number
coincidences~\cite{dirac1938,wesson1980}. Unlike Dirac's original hypothesis,
Eq.~(\ref{eq:main}) is \emph{static}: it relates present-epoch quantities and
carries no implication that $G$ varies with cosmic time. It therefore does not
inherit the observational difficulties that constrain a time-varying $G$.

\section{Discussion: what this does and does not establish}
\label{sec:discussion}

We regard the following distinctions as essential.

\emph{What we report.} A numerical relation among measured quantities,
Eq.~(\ref{eq:main}), which holds to $0.0135\%$ ($0.17\,\sigma$) and which is
the unique sub-$0.1\%$ member of an $83\,160$-candidate family, with an
empirical chance probability below $10^{-4}$.

\emph{What we do not claim.} We do not claim a derivation of
Eq.~(\ref{eq:main}) from an underlying dynamics, nor a mechanism that fixes its
coefficients. The prefactor $4/3$ and the exponent $12$ are not derived here.
We have surveyed the three adjacent traditions in which one might expect the
coefficient $4/3$ to originate---SU(3) flavour symmetry, the geometric reading
of the Koide relation, and the large-number tradition---and in none of them is
$4/3$ produced as a derived quantity. We therefore present the origin of the
coefficient as a genuinely open problem, of a difficulty comparable to the
four-decade-old problem of deriving the Koide relation.

\emph{On the choice of $m_e$.} The reference mass in $\alpha_G$ is
conventional. Its role here is self-consistent rather than fundamental: the
same electron mass appears in the ratio $m_\tau/m_e$ and in $\alpha_G$, so that
both sides of Eq.~(\ref{eq:main}) are expressed in electron units.

\section{Reproducibility}

All quantities reported here are computed from the published constants of
Table~\ref{tab:inputs} by elementary arithmetic. The look-elsewhere enumeration
of Sec.~\ref{sec:lookelsewhere} is a deterministic scan; a script that
regenerates Table~\ref{tab:lookelsewhere} in full, together with the evaluation
of Eq.~(\ref{eq:main}), is available from the author on request and is intended
for public deposit alongside this note.

\section{Conclusion}

We have placed on the record a single, checkable relation,
$(4/3)(m_\tau/m_e)^{12} = \alpha_{\mathrm{EM}}/\alpha_G$, that agrees with
measurement to $0.0135\%$, that is uniquely accurate within a large family of
comparable relations, and whose coefficients admit an exceptional-algebra
reading whose origin remains to be explained. Whether the relation reflects a
deep connection between the lepton mass hierarchy and the gravitational
coupling, or an as-yet-unexplained coincidence, is a question we leave open;
the forthcoming Belle~II tau-mass measurement will test it decisively.

\begin{acknowledgments}
The author thanks S.~Boyko for performing the independent look-elsewhere
analysis of Sec.~\ref{sec:lookelsewhere} and for a critical reading of the
manuscript.
\end{acknowledgments}

\begin{thebibliography}{99}

\bibitem{koide1983} Y.~Koide,
\textit{Phys.\ Lett.\ B} \textbf{120}, 161 (1983).

\bibitem{dirac1938} P.~A.~M.~Dirac,
\textit{Proc.\ R.\ Soc.\ Lond.\ A} \textbf{165}, 199 (1938).

\bibitem{wesson1980} P.~S.~Wesson,
\textit{Gravity, Particles, and Astrophysics} (Reidel, Dordrecht, 1980).

\bibitem{pdg2024} R.~L.~Workman \textit{et al.} (Particle Data Group),
\textit{Prog.\ Theor.\ Exp.\ Phys.} \textbf{2024}, 083C01 (2024).

\bibitem{codata2018} E.~Tiesinga, P.~J.~Mohr, D.~B.~Newell, and B.~N.~Taylor,
\textit{Rev.\ Mod.\ Phys.} \textbf{93}, 025010 (2021) [CODATA~2018].

\bibitem{belleii} W.~Altmannshofer \textit{et al.} (Belle~II),
\textit{Prog.\ Theor.\ Exp.\ Phys.} \textbf{2019}, 123C01 (2019).

\bibitem{bincer1993} A.~M.~Bincer and K.~Riesselmann,
\textit{J.\ Math.\ Phys.} \textbf{34}, 5935 (1993).

\bibitem{sumino2009} Y.~Sumino,
\textit{J.\ High Energy Phys.} \textbf{05} (2009) 075.

\bibitem{singh2021} V.~Bhatta, R.~Mondal, V.~Vaibhav, and T.~P.~Singh,
arXiv:2108.05787 [hep-ph] (2021).

\end{thebibliography}

\end{document}
